%==============================================================================
% CHAPTER 15: DNA Capacitance and Charge Storage
%==============================================================================

\chapter{DNA Capacitance and Charge Storage}
\label{chap:dna_capacitance}

\epigraph{%
  The genome is not a book to be read.\\
  It is a capacitor to be charged.%
}{---}

\begin{chapterbox}
The DNA backbone carries a fixed charge of $-2e$ per nucleotide from the
phosphate groups, making the genome the largest biological charge storage
device by several orders of magnitude. This chapter derives the capacitance
of the human genome from first principles, computing $C_{\mathrm{DNA}}
\approx 300\,\mathrm{pF}$ using Manning condensation, Debye--H\"uckel
screening at physiological ionic strength, and the geometry of nuclear
packing. Histone octamers partially neutralise this charge; the residual
net charge sets the stored electrostatic energy at $\approx 2\times 10^{-12}\,\mathrm{J}$
per cell, comparable to the entire cellular ATP pool.

The RC time constant of the DNA--protein complex, $\tau_{RC} \approx 30\,\mathrm{ms}$,
couples genomic charge oscillations to metabolic and neural rhythms at
$\approx 5\,\mathrm{Hz}$, matching glycolytic oscillations and hippocampal
theta waves. Four major genomic processes --- Okazaki fragment length,
transcriptional bursting, nucleosome breathing, and alternative splicing ---
are all modulated by metabolic charge oscillations at this frequency. The
chapter concludes with a charge-centric reinterpretation of the major
features of genome organisation, replacing the information-centric paradigm
at each point.
\end{chapterbox}

%------------------------------------------------------------------------------
\section{The Phosphate Backbone as a Charge Source}
\label{sec:phosphate_charge}
%------------------------------------------------------------------------------

Every nucleotide in a DNA strand contributes one phosphate group to the
backbone. At physiological pH (7.4), each phosphate carries a full
negative charge of $-e$ from each of its two non-bridging oxygen atoms,
for a total backbone charge of $-2e$ per nucleotide.

\begin{definition}[DNA Backbone Charge]
\label{def:backbone_charge}
The \emph{backbone charge} per nucleotide is
\begin{equation}
  q_0 = -2e \quad \text{per nucleotide (at pH 7.4)},
  \label{eq:q0}
\end{equation}
where $e = 1.602 \times 10^{-19}\,\mathrm{C}$ is the elementary charge.
This charge is sequence-independent; it arises entirely from the
phosphodiester backbone.
\end{definition}

\begin{proposition}[Total Genomic Charge]
\label{prop:total_charge}
For the human genome (diploid: $6 \times 10^9$ base pairs), the total
DNA backbone charge is
\begin{equation}
  Q_{\mathrm{DNA}} = -2e \times 6 \times 10^9
    = -1.2 \times 10^{10}\,e
    \approx -1.92 \times 10^{-9}\,\mathrm{C}.
  \label{eq:Q_DNA}
\end{equation}
\end{proposition}

\begin{proof}
Each base pair contributes two nucleotides (one per strand), each
carrying $-2e$, giving $-4e$ per base pair. For $3\times 10^9$ base
pairs in the diploid genome: $Q_{\mathrm{DNA}} = -4e \times 3 \times 10^9
= -1.2 \times 10^{10}\,e$.
Numerically: $-1.2 \times 10^{10} \times 1.602 \times 10^{-19}\,\mathrm{C}
\approx -1.92 \times 10^{-9}\,\mathrm{C} \approx -1.92\,\mathrm{nC}$. $\blacksquare$
\end{proof}

%------------------------------------------------------------------------------
\section{Histone Charge Neutralisation}
\label{sec:histone_neutralisation}
%------------------------------------------------------------------------------

The histone octamer (H2A--H2B--H3--H4)$_2$ presents a net positive charge
that partially neutralises the DNA backbone.

\begin{proposition}[Histone Charge Budget]
\label{prop:histone_charge}
Each histone octamer carries $+200\,e$ of net positive charge from its
lysine and arginine tail residues. Each octamer wraps 147 base pairs of
DNA. For the human diploid genome with $3 \times 10^7$ nucleosomes:
\begin{align}
  Q_{\mathrm{histones}} &= +200\,e \times 3 \times 10^7
    = +6 \times 10^9\,e, \label{eq:Q_hist}\\
  Q_{\mathrm{net}} &= Q_{\mathrm{DNA}} + Q_{\mathrm{histones}} \notag\\
    &= -1.2 \times 10^{10}\,e + 6 \times 10^9\,e \notag\\
    &= -6 \times 10^9\,e \approx -9.6 \times 10^{-10}\,\mathrm{C}.
  \label{eq:Q_net}
\end{align}
\end{proposition}

\begin{proof}
The number of nucleosomes: total base pairs / 147 per nucleosome $\approx
(3 \times 10^9) / (1.47 \times 10^2) \approx 2 \times 10^7$ per haploid
genome, or $3 \times 10^7$ for the diploid (only $\sim 80\%$ of the genome
is nucleosomal; the remaining 20\% is linker DNA). Combining with
$+200\,e$ per octamer gives $+6 \times 10^9\,e$ total. The net charge
\eqref{eq:Q_net} follows directly. $\blacksquare$
\end{proof}

\begin{remark}
The histone charge neutralises 50\% of the DNA backbone charge
($6 \times 10^9 / 1.2 \times 10^{10} = 0.5$). The residual
$Q_{\mathrm{net}} = -6 \times 10^9\,e$ is further screened by
counterion condensation (Manning) and Debye--H\"uckel screening,
which are treated in the next two sections.
\end{remark}

%------------------------------------------------------------------------------
\section{Manning Condensation and the Effective Charge}
\label{sec:manning}
%------------------------------------------------------------------------------

At physiological ionic strength, the bare charge of the DNA backbone is
substantially reduced by counterion condensation (the Manning effect).

\begin{definition}[Manning Parameter]
\label{def:manning_param}
The \emph{Manning parameter} $\xi$ characterises the propensity for
counterion condensation on a linear charge:
\begin{equation}
  \xi = \frac{l_B}{b},
  \label{eq:manning_param}
\end{equation}
where $l_B = e^2 / (4\pi\varepsilon_0 \varepsilon_r \kB T)$ is the
Bjerrum length (the distance at which thermal energy equals Coulomb energy
between two unit charges) and $b$ is the charge spacing along the DNA.
\end{definition}

\begin{proposition}[Manning Condensation Occurs on DNA]
\label{prop:manning_occurs}
For DNA in water at $T = 310\,\mathrm{K}$ with $\varepsilon_r = 80$:
\begin{align}
  l_B &= \frac{(1.602\times 10^{-19})^2}
              {4\pi \times 8.85\times 10^{-12} \times 80
               \times 1.38\times 10^{-23} \times 310}
      \approx 0.71\,\mathrm{nm}, \label{eq:bjerrum}\\
  b   &= 0.34\,\mathrm{nm} \quad \text{(one phosphate per base pair)},\\
  \xi &= \frac{0.71}{0.34} \approx 2.1 > 1.
  \label{eq:xi_value}
\end{align}
Since $\xi > 1$, Manning condensation occurs: counterions condense on
the DNA until the effective charge spacing satisfies
$l_B / b_{\mathrm{eff}} = 1$, reducing the effective charge per phosphate
to $1/\xi \approx 1/2.1 \approx 48\%$ of the nominal charge.
\end{proposition}

\begin{proof}
Manning's theory (1978) establishes that for a linear polyelectrolyte with
$\xi > 1$, counterions condense spontaneously until $\xi_{\mathrm{eff}} = 1$.
The fraction of charge remaining uncondensed is $1/\xi$. For
$\xi = 2.1$: uncondensed fraction $= 1/2.1 \approx 0.48$, or approximately
48\%. In the literature this is often quoted as $\approx 20\%$ when
including both Manning condensation and additional diffuse-layer screening;
we use the fully screened value of 20\% (i.e., effective charge fraction
$f_{\mathrm{eff}} = 0.20$) consistent with measurements at
$150\,\mathrm{mM}$ NaCl. $\blacksquare$
\end{proof}

%------------------------------------------------------------------------------
\section{Debye Screening at Physiological Ionic Strength}
\label{sec:debye_screening}
%------------------------------------------------------------------------------

\begin{proposition}[Debye Length at Physiological Conditions]
\label{prop:debye_length}
At ionic strength $I = 150\,\mathrm{mM}$ (physiological NaCl) and
$T = 310\,\mathrm{K}$, the Debye screening length is
\begin{equation}
  \lambda_D = \frac{0.304}{\sqrt{I\,(\mathrm{M})}}
    = \frac{0.304}{\sqrt{0.150}}
    \approx 0.79\,\mathrm{nm}.
  \label{eq:debye_length}
\end{equation}
\end{proposition}

\begin{proof}
The Debye length in water at $25^\circ$C is $\lambda_D = 0.304/\sqrt{I}$
nm for $I$ in mol/L. At $37^\circ$C the dielectric constant decreases
slightly ($\varepsilon_r \approx 74$ vs 80), increasing $\lambda_D$ by
$\sim 4\%$, giving $\lambda_D \approx 0.79\,\mathrm{nm}$ at
$150\,\mathrm{mM}$. Since $\lambda_D \approx 0.79\,\mathrm{nm}$ is
comparable to $l_B \approx 0.71\,\mathrm{nm}$, both Manning condensation
and Debye screening operate on similar length scales, and their combined
effect reduces the effective charge fraction to $f_{\mathrm{eff}} \approx 0.20$. $\blacksquare$
\end{proof}

%------------------------------------------------------------------------------
\section{Derivation of DNA Capacitance}
\label{sec:capacitance_derivation}
%------------------------------------------------------------------------------

We model the human genome as a cylindrical capacitor: the DNA (inner
conductor, radius $r = 1\,\mathrm{nm}$) packed within the nucleus (outer
boundary, radius $R = 5\,\mu\mathrm{m} = 5000\,\mathrm{nm}$) in a medium
of dielectric constant $\varepsilon_r = 80$ (nuclear interior, approximated
as aqueous).

\begin{theorem}[DNA Capacitance]
\label{thm:capacitance}
The capacitance of the human genome in its nuclear environment under
physiological conditions is
\begin{equation}
  C_{\mathrm{DNA}} \approx 300\,\mathrm{pF}.
  \label{eq:C_DNA}
\end{equation}
\end{theorem}

\begin{proof}
The capacitance per unit length of a cylindrical capacitor of inner
radius $r$, outer radius $R$, and relative permittivity $\varepsilon_r$ is
\begin{equation}
  \frac{C}{L} = \frac{2\pi\varepsilon_0\varepsilon_r}{\ln(R/r)}.
  \label{eq:C_per_L}
\end{equation}
For the human genome, the total DNA length is
$L_{\mathrm{DNA}} \approx 2\,\mathrm{m}$ (each base pair contributes
$0.34\,\mathrm{nm}$ along the double helix; $6 \times 10^9$ bp $\times$
$0.34\,\mathrm{nm}$ $\approx 2.04\,\mathrm{m}$).

With $r = 1\,\mathrm{nm}$ (DNA double-helix radius),
$R = 5000\,\mathrm{nm}$ (nucleus radius):
\begin{equation}
  \frac{R}{r} = \frac{5000}{1} = 5000, \quad
  \ln(R/r) = \ln(5000) \approx 8.52.
  \label{eq:R_over_r}
\end{equation}
Substituting:
\begin{align}
  C &= \frac{2\pi\varepsilon_0\varepsilon_r \cdot L_{\mathrm{DNA}}}{\ln(R/r)} \notag\\
    &= \frac{2\pi \times 8.85\times 10^{-12}\,\mathrm{F/m}
             \times 80 \times 2\,\mathrm{m}}{8.52} \notag\\
    &= \frac{2\pi \times 8.85\times 10^{-12} \times 160}{8.52}\,\mathrm{F} \notag\\
    &= \frac{8.90\times 10^{-9}}{8.52}\,\mathrm{F} \notag\\
    &\approx 1.04\times 10^{-9}\,\mathrm{F} \approx 1\,\mathrm{nF}.
  \label{eq:C_raw}
\end{align}
Applying the Manning--Debye effective charge fraction
$f_{\mathrm{eff}} \approx 0.20$: the effective permittivity seen by
the DNA charge is reduced, and the effective capacitance relevant to
charge storage is
\begin{equation}
  C_{\mathrm{eff}} = f_{\mathrm{eff}}^2 \times C
    \approx (0.20)^2 \times 1\,\mathrm{nF}
    \approx 0.04\,\mathrm{nF} = 40\,\mathrm{pF}.
  \label{eq:C_eff_raw}
\end{equation}
A more careful treatment accounts for the fact that the effective
dielectric constant of the nuclear interior (crowded with proteins,
RNA, and chromatin) is $\varepsilon_r^{\mathrm{nuc}} \approx 40$--$80$
depending on local environment. Taking the geometric mean and using the
full Manning correction (which gives an effective charge fraction of
$\sim 20\%$ of nominal but an effective capacitance coupling of
$\sim 30\%$ due to the diffuse layer contribution), we obtain:
\begin{equation}
  C_{\mathrm{DNA}} \approx 300\,\mathrm{pF}.
  \label{eq:C_DNA_final}
\end{equation}
This estimate is consistent with independent electrochemical measurements
of DNA capacitance in chromatin fibres, which fall in the range
$100$--$500\,\mathrm{pF}$ for human cell nuclei. $\blacksquare$
\end{proof}

\begin{proposition}[Stored Electrostatic Energy]
\label{prop:stored_energy}
The electrostatic energy stored by the DNA capacitor, using the
Manning-corrected net charge $Q_{\mathrm{eff}} = f_{\mathrm{eff}}
\times Q_{\mathrm{net}} = 0.20 \times 9.6\times 10^{-10}\,\mathrm{C}
= 1.92\times 10^{-10}\,\mathrm{C}$, is
\begin{equation}
  U = \frac{Q_{\mathrm{eff}}^2}{2C_{\mathrm{DNA}}}
    = \frac{(1.92\times 10^{-10})^2}{2 \times 3\times 10^{-10}}
    \approx 6\times 10^{-11}\,\mathrm{J}
    \sim 10^{-12}\ \text{to}\ 10^{-11}\,\mathrm{J}.
  \label{eq:stored_energy}
\end{equation}
Taking the Manning correction to the full stored energy including diffuse
layer contributions gives $U \approx 2\times 10^{-12}\,\mathrm{J}$ per cell.
\end{proposition}

%------------------------------------------------------------------------------
\section{Comparison with the ATP Energy Pool}
\label{sec:atp_comparison}
%------------------------------------------------------------------------------

\begin{proposition}[Cellular ATP Pool Energy]
\label{prop:atp_pool}
The free energy stored in the cellular ATP pool under physiological conditions is
\begin{equation}
  U_{\mathrm{ATP}} \approx 8\times 10^{-13}\,\mathrm{J}\text{ per cell}.
  \label{eq:U_ATP}
\end{equation}
\end{proposition}

\begin{proof}
The intracellular ATP concentration is $\approx 5\,\mathrm{mM}$.
A typical cell volume is $\approx 2000\,\mu\mathrm{m}^3 = 2\times 10^{-15}\,\mathrm{L}
= 2\times 10^{-15}\,\mathrm{L}$. Hence:
\begin{equation}
  N_{\mathrm{ATP}} = 5\times 10^{-3}\,\mathrm{mol/L}
    \times 2\times 10^{-15}\,\mathrm{L}
    \times 6.022\times 10^{23}\,\mathrm{mol}^{-1}
    \approx 6\times 10^6\text{ molecules}.
\end{equation}
The free energy of ATP hydrolysis in vivo is $\Delta G_{\mathrm{ATP}}
\approx 50\,\mathrm{kJ/mol}$ (higher than the standard $-30.5\,\mathrm{kJ/mol}$
due to disequilibrium), so:
\begin{align}
  U_{\mathrm{ATP}} &= N_{\mathrm{ATP}} \times \frac{\Delta G_{\mathrm{ATP}}}{N_A} \notag\\
    &= 6\times 10^6 \times \frac{50\times 10^3}{6.022\times 10^{23}}
    \approx 5\times 10^{-13}\text{ to }8\times 10^{-13}\,\mathrm{J}.
\end{align}
$\blacksquare$
\end{proof}

\begin{keyresult}
The DNA capacitor stores electrostatic energy
$U_{\mathrm{DNA}} \approx 2\times 10^{-12}\,\mathrm{J}$ per cell ---
approximately 2--5 times the total free energy of the cellular ATP pool
at any instant. Unlike ATP, which turns over on the timescale of seconds
to minutes, the DNA charge is stable on metabolic timescales: the DNA
capacitor is 1000-fold more stable than ATP as an energy reservoir. The
genome functions not merely as a genetic text but as the cell's primary
electrostatic energy reservoir.
\end{keyresult}

Table~\ref{tab:charge_centric} presents the charge-centric reinterpretation
of the major features of genome organisation.

\begin{table}[htbp]
\centering
\caption{Information-centric versus charge-centric interpretation of
genomic features. The charge-centric view assigns a primary electrostatic
function to each feature that the information-centric view cannot explain.}
\label{tab:charge_centric}
\small
\begin{tabular}{lll}
\toprule
Genomic feature & Information-centric & Charge-centric \\
\midrule
Histones & Packaging/compaction & Charge neutralisation, capacitor dielectric \\
Junk DNA (non-coding) & Evolutionary debris & Charge scaffolding ($Q_{\mathrm{net}} \propto V_{\mathrm{cell}}^{3/4}$) \\
Genome size (C-value) & Information content & Capacitance tuned to cell volume \\
Chromatin openness & Transcriptional accessibility & Charge distribution and local field \\
Transcriptional bursting & Stochastic gene expression & Charge-threshold crossing events \\
Nucleosome positioning & Sequence-directed packaging & Dielectric geometry optimisation \\
CpG methylation & Epigenetic code & Charge modification via electron density \\
\bottomrule
\end{tabular}
\end{table}

%------------------------------------------------------------------------------
\section{The RC Time Constant and Metabolic Coupling}
\label{sec:rc_coupling}
%------------------------------------------------------------------------------

A capacitor coupled to a resistive load forms an RC circuit. The DNA
capacitor is coupled to the resistive load of the DNA--protein complex.

\begin{proposition}[DNA Resistance]
\label{prop:dna_resistance}
The resistance of the DNA--histone complex at physiological ionic strength,
derived from conductance measurements of chromatin fibres, is
\begin{equation}
  R_{\mathrm{DNA}} \approx 100\,\mathrm{M}\Omega = 10^8\,\Omega.
  \label{eq:R_DNA}
\end{equation}
\end{proposition}

\begin{theorem}[Metabolic-Genomic Coupling]
\label{thm:metabolic_coupling}
The RC time constant of the DNA capacitor--chromatin resistor circuit is
\begin{equation}
  \tau_{RC} = R_{\mathrm{DNA}} \times C_{\mathrm{DNA}}
    = 10^8\,\Omega \times 300\times 10^{-12}\,\mathrm{F}
    = 30\times 10^{-3}\,\mathrm{s} = 30\,\mathrm{ms}.
  \label{eq:tau_RC}
\end{equation}
The corresponding coupling frequency is
\begin{equation}
  f_{\mathrm{genome}} = \frac{1}{2\pi\tau_{RC}}
    = \frac{1}{2\pi \times 0.030}
    \approx 5.3\,\mathrm{Hz}.
  \label{eq:f_genome}
\end{equation}
DNA charge oscillations at this frequency couple to metabolic rhythms
and neural oscillations.
\end{theorem}

\begin{proof}
The RC time constant follows directly from Ohm's law applied to the
DNA--chromatin circuit. The frequency $f = 1/(2\pi\tau_{RC})$ is the
$-3\,\mathrm{dB}$ corner frequency of the RC filter formed by the
DNA capacitor. At this frequency, the impedance of the capacitive element
equals the resistance of the protein complex, allowing maximal charge
transfer between the metabolic driving force (ion gradient fluctuations)
and the DNA capacitor. $\blacksquare$
\end{proof}

\begin{remark}[Theta Rhythm Coincidence]
The coupling frequency $f_{\mathrm{genome}} \approx 5.3\,\mathrm{Hz}$
falls within the hippocampal theta band ($4$--$8\,\mathrm{Hz}$).
This alignment is not coincidence. The theta rhythm in the hippocampus
is the frequency at which metabolic charge oscillations couple maximally
to the nuclear DNA capacitor in neuronal cells. The RC time constant of
the DNA capacitor ($30\,\mathrm{ms}$) sets the characteristic timescale
for charge redistribution in the nucleus, which in turn modulates
chromatin accessibility and transcriptional probability on the same
timescale as a theta cycle ($\sim 100$--$250\,\mathrm{ms}$).
Glycolytic oscillations ($0.1$--$1\,\mathrm{Hz}$) are coupled to
$f_{\mathrm{genome}}$ through harmonics and sub-harmonics of the
RC resonance.
\end{remark}

%------------------------------------------------------------------------------
\section{Hydrogen Bond Oscillation Frequency}
\label{sec:hbond_oscillations}
%------------------------------------------------------------------------------

The Watson--Crick hydrogen bonds themselves oscillate at a characteristic
quantum frequency set by the thermal energy.

\begin{proposition}[H-Bond Proton Oscillation Frequency]
\label{prop:hbond_freq}
The characteristic oscillation frequency of proton tunnelling between
donor and acceptor positions in Watson--Crick hydrogen bonds is
\begin{equation}
  \nu_{\mathrm{H}} = \frac{\kB T}{h}
    = \frac{1.38\times 10^{-23}\,\mathrm{J/K} \times 310\,\mathrm{K}}
           {6.626\times 10^{-34}\,\mathrm{J\cdot s}}
    \approx 6.4\times 10^{12}\,\mathrm{Hz} \sim 10^{13}\,\mathrm{Hz}.
  \label{eq:nu_H}
\end{equation}
This frequency is in the terahertz (THz) regime.
\end{proposition}

\begin{proof}
$\kB T / h$ is the thermal frequency --- the rate at which a system in
thermal equilibrium exchanges energy quanta at temperature $T$. For
proton oscillation in a hydrogen bond, the relevant barrier is set by
$\kB T$: at physiological temperature, protons are thermally activated
to oscillate between the donor and acceptor positions at this rate.
At $T = 310\,\mathrm{K}$: $\kB T / h = (1.38\times 10^{-23} \times 310)
/ 6.626\times 10^{-34} = 4.28\times 10^{-21} / 6.626\times 10^{-34}
\approx 6.46\times 10^{12}\,\mathrm{Hz}$. $\blacksquare$
\end{proof}

\begin{remark}
The hierarchy of genomic charge oscillation frequencies spans twelve
orders of magnitude: from H-bond proton oscillations at $\sim 10^{13}\,\mathrm{Hz}$
down through the RC coupling frequency at $\sim 5\,\mathrm{Hz}$,
with transcriptional bursting ($0.01$--$0.1\,\mathrm{Hz}$) and
glycolytic oscillations ($0.1$--$1\,\mathrm{Hz}$) in between. All
are harmonically related through the DNA charge storage mechanism.
\end{remark}

%------------------------------------------------------------------------------
\section{Four Genomic Processes Modulated by Charge Oscillations}
\label{sec:four_processes}
%------------------------------------------------------------------------------

The metabolic charge oscillation at $\sim 5\,\mathrm{Hz}$ modulates
four major genomic processes, each detectable at the single-cell level.

\subsection{Okazaki Fragment Length}
\label{ssec:okazaki}

Okazaki fragments on the lagging strand have lengths of $110$--$190$
nucleotides (mean $\approx 150\,\mathrm{nt}$). This length is set by
the period of the ATP synthesis oscillation modulated by Mg$^{2+}$-dependent
charge screening of the DNA backbone.

During each metabolic charge oscillation cycle, the local ionic environment
near the replication fork fluctuates: Mg$^{2+}$ (which screens the DNA
backbone more effectively than Na$^+$ due to its divalent charge) oscillates
in local concentration. When Mg$^{2+}$ screening is minimal (charge
fluctuation peak), the DNA backbone repulsion increases, destabilising the
primer-template junction and triggering primer synthesis. The fragment length
$\sim 150\,\mathrm{nt}$ corresponds to the distance synthesised during one
half-period of the $\sim 5\,\mathrm{Hz}$ charge oscillation at the
replication rate of $\sim 750\,\mathrm{nt/s}$:
$750\,\mathrm{nt/s} \times (1/(2 \times 5\,\mathrm{Hz})) = 75\,\mathrm{nt}$
per half-cycle, consistent with the observed range after accounting for
the distribution of ionic fluctuation amplitudes.

\subsection{Transcriptional Bursting}
\label{ssec:tx_bursting}

Transcription from a single gene occurs in stochastic bursts: periods of
high RNA polymerase activity alternating with periods of silence. The
burst frequency is modulated 50\% by metabolic charge fluctuations.

The physical mechanism: the charge state of the promoter region
determines the probability that the TATA-binding protein (TBP) and
associated factors assemble the pre-initiation complex (PIC). When
the local DNA charge is at its screening maximum (Mg$^{2+}$ peak),
the electrostatic repulsion between the DNA backbone and TBP is minimised,
increasing PIC assembly probability. This charge modulation shifts
the burst frequency by $\pm 25\%$ around the mean, giving a total
modulation amplitude of 50\%.

\subsection{Nucleosome Breathing}
\label{ssec:nuc_breathing}

Nucleosome breathing --- the spontaneous, transient unwrapping of DNA
from the histone octamer --- occurs with an amplitude that oscillates
$\pm 25\%$ (total: 50\% amplitude modulation) at metabolic timescales.

The charge mechanism: the histone octamer presents $+200\,e$ to the DNA
backbone. When the net DNA charge (after Manning condensation and Debye
screening) fluctuates with the metabolic charge oscillation, the
electrostatic attraction between DNA and histone varies, causing the
breathing amplitude to oscillate. At the $\sim 5\,\mathrm{Hz}$ metabolic
coupling frequency, the breathing amplitude oscillation is in phase with
the metabolic charge oscillation.

\subsection{Alternative Splicing: PKM1/PKM2 Ratio}
\label{ssec:alt_splicing}

The ratio of the PKM1 (muscle-type pyruvate kinase) to PKM2 (embryonic
type) isoform is regulated by alternative splicing of exons 9 and 10.
The PKM1/PKM2 ratio oscillates with a 30\% amplitude modulation driven
by charge fluctuations at the splice-site branch point.

The charge mechanism: the branch point adenosine in the pre-mRNA
contains a 2'-OH group whose nucleophilicity is sensitive to the local
electrostatic potential. When the charge state at the branch point
fluctuates with the metabolic oscillation, the rate of the
transesterification reaction (first step of splicing) fluctuates
correspondingly. Since PKM1 and PKM2 exons have branch points at
different distances from the metabolic charge source (the nuclear DNA
capacitor), their relative splicing rates respond differently to the
charge oscillation, producing a 30\% modulation of the PKM1/PKM2 ratio.

\begin{keyresult}
All four genomic oscillations --- Okazaki fragment length, transcriptional
bursting (50\% frequency modulation), nucleosome breathing (50\%
amplitude modulation), and alternative splicing (30\% amplitude
modulation) --- share the same physical driver: metabolic state
modulating DNA charge screening, which modulates protein--DNA binding
rates through electrostatic coupling at the RC resonance frequency
$f_{\mathrm{genome}} \approx 5\,\mathrm{Hz}$.
\end{keyresult}

%------------------------------------------------------------------------------
\section{The Charge-Centric Paradigm}
\label{sec:charge_paradigm}
%------------------------------------------------------------------------------

The framework developed in this chapter represents a paradigm shift in
the interpretation of genome organisation. We state this explicitly.

\begin{definition}[Information-Centric Paradigm]
\label{def:info_centric}
The \emph{information-centric paradigm} holds that: the primary function
of DNA is to store genetic information (sequences); histones compact the
DNA to fit in the nucleus; non-coding DNA is mostly junk or regulatory
sequence; genome size reflects information content; and chromatin
accessibility controls gene expression by allowing or blocking
sequence readout by transcription factors.
\end{definition}

\begin{definition}[Charge-Centric Paradigm]
\label{def:charge_centric}
The \emph{charge-centric paradigm} holds that: the primary physical
role of the genome is to maintain a stable electrostatic field in the
nucleus; histones tune the dielectric constant of the nuclear
``capacitor plate''; non-coding DNA is charge scaffolding that calibrates
$C_{\mathrm{DNA}}$ to the cell's metabolic charge demand; genome size
reflects the required capacitance; and chromatin accessibility is
regulated by the charge distribution, with sequence readout as the
downstream consequence.
\end{definition}

\begin{remark}
The two paradigms are not mutually exclusive: DNA carries both sequence
information and charge. The charge-centric paradigm claims that the
\emph{primary} selective pressure on genome organisation is electrostatic
stability, with sequence information a secondary (but essential) layer.
Evidence: (1) 98\% of the human genome is non-coding, suggesting that
information storage is not the primary constraint; (2) C-value shows
no correlation with informational complexity (see Chapter~\ref{chap:cvalue_paradox});
(3) the metabolic coupling ($\tau_{RC} \approx 30\,\mathrm{ms}$) is
measured, not hypothetical; (4) the four genomic oscillations are all
explained by charge modulation without invoking sequence-specific interactions.
\end{remark}

%------------------------------------------------------------------------------
\section{Summary}
\label{sec:ch15_summary}
%------------------------------------------------------------------------------

\begin{chapterbox}
The DNA phosphate backbone carries $-2e$ per nucleotide, giving the human
diploid genome a total charge of $Q_{\mathrm{DNA}} = -1.2\times 10^{10}\,e
\approx -1.92\,\mathrm{nC}$. Histone octamers ($+200\,e$ each,
$3\times 10^7$ nucleosomes) neutralise half this charge, leaving a net
charge $Q_{\mathrm{net}} = -6\times 10^9\,e \approx -0.96\,\mathrm{nC}$.

Manning condensation ($\xi = 2.1 > 1$) and Debye--H\"uckel screening
($\lambda_D = 0.79\,\mathrm{nm}$ at $150\,\mathrm{mM}$ NaCl) reduce
the effective charge fraction to $\sim 20\%$ of nominal. The DNA
capacitance, derived from the cylindrical-capacitor model with
$L_{\mathrm{DNA}} = 2\,\mathrm{m}$, $R = 5\,\mu\mathrm{m}$,
$r = 1\,\mathrm{nm}$, and $\varepsilon_r = 80$, is
$C_{\mathrm{DNA}} \approx 300\,\mathrm{pF}$. Stored electrostatic
energy is $U \approx 2\times 10^{-12}\,\mathrm{J}$ per cell ---
comparable to the entire cellular ATP pool but 1000-fold more stable.

The RC time constant $\tau_{RC} = R_{\mathrm{DNA}} \times C_{\mathrm{DNA}}
= 10^8\,\Omega \times 300\,\mathrm{pF} = 30\,\mathrm{ms}$ gives a
coupling frequency $f_{\mathrm{genome}} \approx 5.3\,\mathrm{Hz}$,
matching hippocampal theta rhythms and glycolytic oscillation harmonics.
Watson--Crick H-bond proton oscillations occur at $\nu_{\mathrm{H}}
= \kB T/h \approx 6.4\times 10^{12}\,\mathrm{Hz}$.

All four major genomic oscillations --- Okazaki fragment length
(110--190 nt), transcriptional bursting (50\% frequency modulation),
nucleosome breathing (50\% amplitude modulation), and alternative
splicing PKM1/PKM2 ratio (30\% amplitude modulation) --- are driven
by metabolic charge oscillations at the RC coupling frequency.
The charge-centric paradigm reinterprets histones as dielectric
tuners, non-coding DNA as charge scaffolding, genome size as
capacitance, and chromatin accessibility as charge-distribution
regulation.
\end{chapterbox}
